% ----------------------------------
% In the preliminaries should explain
%   - Probability space
%       - \sigma-algebra
%   - Stochastic processes
%       - Martingales
%       ? Markov property
%   - Brownian motion
% ----------------------------------

\section{Preliminary notions}

\subsection{Probability theory}

The foundation of modern probability theory is based on measure theory.

\begin{definition}[$\sigma$-algebra]
    A \textbf{$\sigma$-algebra} $\F$ on a set $\Omega$ is a collection of subsets of $\Omega$ that
    \begin{itemize}
        \item $\Omega \in \F$
        \item if $A \in \F$ then $X \setminus A \in \F$
        \item if $A_1, A_2, \ldots \in \F$ then $\bigcup_{i = 1}^{\infty} A_i \in \F$
    \end{itemize}
\end{definition}

\begin{definition}[Measure]
    Let $\Omega$ be a set and $\F$ a $\sigma$-algebra on $\Omega$.
    A function $\mu: \F \rightarrow [0, \infty]$ is \textbf{measure} if satisfies
    \begin{enumerate}
        \item $\mu(\emptyset) = 0$
        \item $\mu(A) \geq 0$ for all $A \in \F$
        \item if $\{ A_i \}_{i = 1}^{\infty} \in \F$ and $A_i \cap A_j = \emptyset$ for all $i \neq j$ then
            \begin{equation*}
                \mu \left( \bigcup_{i = 1}^{\infty} A_i \right) 
                = \sum_{i = 1}^{\infty} \mu \left( A_i \right)
            \end{equation*}
    \end{enumerate}
\end{definition}

If $\mu$ is a measure on $\F$ such that $\mu(\Omega) = 1$,
then $\mu$ is called a \textbf{probability measure}.

\begin{definition}[Probability space]
    A probability space is a triple $(\Omega, \F, P)$ where
    \begin{itemize}
        \item $\Omega$ is the set of all possible outcomes
        \item $\F$ is a $\sigma$-algebra on $\Omega$
        \item $P$ is a probability measure.
    \end{itemize}
\end{definition}

For a more details in probability theory and stochastic processes, see \cite{kallenberg_probability} or \cite{klenke_probability}.

\subsection{Stochastic processes}

\begin{definition}[Stochastic processes]
    A stochastic process is an indexed family of random variables
    \begin{equation*}
        \{ X(t, \omega) \}_{t \in T},
    \end{equation*}
    on a probability space $(\Omega, \F, P)$ assuming values in $\R^n$.
\end{definition}

If the set $T$ is countable then the stochastic process is said to be \textit{discrete},
if $T$ is some interval of $\R$ the is \textit{continuous}.

An intuitive way to understand stochastic processes is to consider what is $X(t, \omega)$ whan $t$ or $\omega$ is fixed.
If we fix $t \in T$ then the map
\begin{equation*}
    \omega \rightarrow X_t(\omega),
\end{equation*}
for $\omega \in \Omega$ is a random variable.
If $\omega \in \Omega$ is fixed then the map
\begin{equation*}
    t \rightarrow X_\omega(t),
\end{equation*}
for $t \in T$ is a deterministic function, and is called \textit{path} of $X(t, \omega)$.


\begin{definition}[Filtration]
A \textbf{filtration} is a family of $\sigma$-algebras $\{ \F_t \}_{t \in T}$ indexed by $T$,
such that $\F_s \subseteq \F_t$ for all $s \leq t$, where $s, t \in T$.
\end{definition}

The idea is that a filtration $\F_t$ represents the information available form the past of an stochastic process up to time $t$.
Thus we define a \textbf{filtered probability space} as a probability space with a filtration $(\Omega, \F, \{ \F_t \}_{t \in T}, P)$.

\begin{definition}[Adapted process]
    Let $X_t$ be a stochastic process on a filtered probability space $(\Omega, \F, \{ \F_t \}_{t \in T}, P)$.
    The process is \textbf{adapted} to the filtration $\{ \F_t \}_{t \in T}$ 
    if, for every $t \in T$,
    the random variable $X_t$ is $F_t$-measurable.
\end{definition}

The idea of an adapted process is key when modeling using stochastic processes,
as it ensures that the process is at time $t$ only depends on the information available up to that time,
but not on future information.

\subsubsection{Martingales}

A martingale is a fundamental concept in the theory of stochastic processes,
formalizes the idea of a "fair game",
where given all the information up to now,
the best prediction for the future is simply the present value.

\begin{definition}[Martingale]
    A stochastic process ${X}_{t}$ where $t \in T$ is called a \textbf{martingale}
    with respect to a filtration ${F}_{t}$ if it satisfies the following conditions:
    \begin{enumerate}
        \item $E[|X_t|] < \infty$, for each $t \in T$ 
        \item $X_t$ is measurable with respect to $\F_t$ for each $t \in T$
        \item If $X_t$ is discrete
            \begin{equation*}
                E[X_{t + 1} | F_t] = X_t \quad \forall t \in \N,
            \end{equation*}
            if $X_t$ is continuous
            \begin{equation*}
                E[X_t | F_s] = X_s \quad \forall s \leq t.
            \end{equation*}
    \end{enumerate}
\end{definition}

% Local martingales ??

% Markov property ??

% Stopping time ??

\subsubsection{Brownian Motion}

Brownian motion (also called Wiener process) is the fundamental stochastic process underlying continuous-time financial modeling.
Formally, a Brownian motion $\{ W_t \}_{(t \geq 0)}$ is defined by the following properties:

\begin{enumerate}
    \item $W_0 = 0$

    \item The map $t \rightarrow W_t(\omega)$ is almost surely continuous

    \item $W_t$ has independent increments,
        i.e. For $t_1 < t_2 < t_3 < t_4$, $W_{t_2} - W_{t_2}$ and $W_{t_4} - W_{t_4}$ are independent random variables

    \item $W_t$ has gaussian increments,
        i.e. $W_t - W_s \sim \mathcal{N}(0, t-s)$ for $0 \leq s \leq t$
\end{enumerate}
where $\mathcal{N}(\mu, \sigma)$ is the normal distribution with mean $\mu$ and variance $\sigma$.

The Brownian motion has some key features that will need afterwards,
for example it can be shown that the paths $t \rightarrow W_t(\omega)$ although are almost surely continuous,
they are nowhere differentiable.

Also the it can be shown that the Wiener process is an almost surely continuous martingale,
and its quadratic variation is $\lbrace W_t, W_t \rbrace = t$.
This two properties along $W_0 = 0$ form an equivalent definition of the Wiener process,
called the Levy's characterization.

In \cite{morters_peres_brownian_motion} there are profs of this properties and the equivalence between the tow definitions,
also there is a prof of the existence of the Brownian motion.
